% GENERATED FILE. Edit canonical lessons, then run npm run generate:books.
% canonical-source-hash: fnv1a64:daa818cddc3f34e0
% canonical-lessons: JA-C299-ousama, JA-C299-ohimesama, JA-C299-samurai, JA-C299-mamoru, JA-C299-tonosama

\chapter{The King's Daughter}
\label{ch:ja-the-king-s-daughter}
\hlchaptermodality{299}

\begin{chapteropening}
\textbf{By the end of this chapter:} \emph{I can say who ruled and who protected them.}

Complete the last lesson of chapter 299: I can say who ruled and who protected them.
\end{chapteropening}

\section[\texorpdfstring{\ja{おうさま} (ōsama)}{おうさま (ōsama)}]{\ja{おうさま} (ōsama) — a king}

\label{lesson:JA-C299-ousama}



\begin{quote}
Before the new one: say the Japanese for a village, then the Japanese for a hard worker.
\end{quote}

\subsection*{You'll want to know: \ja{おうさま}}
\textbf{\ja{おうさま}} — \emph{ōsama} — \textquotedblleft{}a king\textquotedblright{}.

\begin{quote}
\textbf{\ja{むかし、その} \ja{くにに} \ja{おうさまが} \ja{いました。}} — \emph{mukashi, sono kuni ni ōsama ga imashita} — \textquotedblleft{}Long ago that country had a king.\textquotedblright{}
\end{quote}

\begin{cousinweb}
\textbf{\ja{おうさま}} (\emph{ōsama}) is a mixed word, native and Sino-Japanese: \emph{ō}, \textquotedblleft{}king\textquotedblright{}, is a Chinese character read with its on'yomi, a reading borrowed from Chinese (Middle Chinese \emph{hjwang}), and \emph{sama}, \textquotedblleft{}polite suffix\textquotedblright{}, is native Japanese (JMdict, KANJIDIC2; Middle Chinese in Baxter's transcription).
\end{cousinweb}

\begin{grammarlens}[title={What you've built}]
One more word for telling what happened.
\end{grammarlens}

\subsection*{Your turn}
\begin{itemize}
\raggedright
  \item \emph{Say it:} \emph{ōsama}
  \item \emph{Say it:} \emph{ōsama}, once more
  \item \emph{Say it:} \emph{mukashi, sono kuni ni ōsama ga imashita} — Long ago that country had a king.
  \item \emph{Recall:} say the Japanese for a village, then the Japanese for a hard worker, then say \emph{ōsama} again
\end{itemize}

\begin{tcolorbox}[breakable,colback=teal!4,colframe=teal!35!black,title={Before you move on}]
What does \ja{おうさま} mean? (\textquotedblleft{}A king\textquotedblright{}.) Say it once more.
\end{tcolorbox}
\section[\texorpdfstring{\ja{おひめさま} (ohimesama)}{おひめさま (ohimesama)}]{\ja{おひめさま} (ohimesama) — a princess}

\label{lesson:JA-C299-ohimesama}



\begin{quote}
Before the new one: say the Japanese for a hard worker, then the Japanese for a king.
\end{quote}

\subsection*{You'll want to know: \ja{おひめさま}}
\textbf{\ja{おひめさま}} — \emph{ohimesama} — \textquotedblleft{}a princess\textquotedblright{}.

\begin{quote}
\textbf{\ja{おうさまには} \ja{うつくしい} \ja{おひめさまが} \ja{いました。}} — \emph{ōsama ni wa utsukushii ohimesama ga imashita} — \textquotedblleft{}The king had a beautiful daughter.\textquotedblright{}
\end{quote}

\begin{cousinweb}
\textbf{\ja{おひめさま}} (\emph{ohimesama}) is a native Japanese word (wago), inherited rather than borrowed from Chinese, built from \emph{o-}, \textquotedblleft{}the polite prefix\textquotedblright{}, \emph{hime}, \textquotedblleft{}princess\textquotedblright{}, and \emph{sama}, \textquotedblleft{}the honorific\textquotedblright{} (JMdict, KANJIDIC2).
\end{cousinweb}

\begin{grammarlens}[title={What you've built}]
One more word for telling what happened.
\end{grammarlens}

\subsection*{Your turn}
\begin{itemize}
\raggedright
  \item \emph{Say it:} \emph{ohimesama}
  \item \emph{Say it:} \emph{ohimesama}, once more
  \item \emph{Say it:} \emph{ōsama ni wa utsukushii ohimesama ga imashita} — The king had a beautiful daughter.
  \item \emph{Recall:} say the Japanese for a hard worker, then the Japanese for a king, then say \emph{ohimesama} again
\end{itemize}

\begin{tcolorbox}[breakable,colback=teal!4,colframe=teal!35!black,title={Before you move on}]
What does \ja{おひめさま} mean? (\textquotedblleft{}A princess\textquotedblright{}.) Say it once more.
\end{tcolorbox}
\section[\texorpdfstring{\ja{さむらい} (samurai)}{さむらい (samurai)}]{\ja{さむらい} (samurai) — a samurai; a warrior}

\label{lesson:JA-C299-samurai}



\begin{quote}
Before the new one: say the Japanese for a king, then the Japanese for a princess.
\end{quote}

\subsection*{You'll want to know: \ja{さむらい}}
\textbf{\ja{さむらい}} — \emph{samurai} — \textquotedblleft{}a samurai; a warrior\textquotedblright{}.

\begin{quote}
\textbf{\ja{つよい} \ja{さむらいが} \ja{むらに} \ja{きました。}} — \emph{tsuyoi samurai ga mura ni kimashita} — \textquotedblleft{}A strong samurai came to the village.\textquotedblright{}
\end{quote}

\begin{cousinweb}
\textbf{\ja{さむらい}} (\emph{samurai}) is a native Japanese word (wago), inherited rather than borrowed from Chinese; its kanji was attached later as a kun'yomi (JMdict).
\end{cousinweb}

\begin{grammarlens}[title={What you've built}]
One more word for telling what happened.
\end{grammarlens}

\subsection*{Your turn}
\begin{itemize}
\raggedright
  \item \emph{Say it:} \emph{samurai}
  \item \emph{Say it:} \emph{samurai}, once more
  \item \emph{Say it:} \emph{tsuyoi samurai ga mura ni kimashita} — A strong samurai came to the village.
  \item \emph{Recall:} say the Japanese for a king, then the Japanese for a princess, then say \emph{samurai} again
\end{itemize}

\begin{tcolorbox}[breakable,colback=teal!4,colframe=teal!35!black,title={Before you move on}]
What does \ja{さむらい} mean? (\textquotedblleft{}A samurai; a warrior\textquotedblright{}.) Say it once more.
\end{tcolorbox}
\section[\texorpdfstring{\ja{まもる} (mamoru)}{まもる (mamoru)}]{\ja{まもる} (mamoru) — to protect; to guard}

\label{lesson:JA-C299-mamoru}



\begin{quote}
Before the new one: say the Japanese for a princess, then the Japanese for a samurai.
\end{quote}

\subsection*{You'll want to know: \ja{まもる}}
\textbf{\ja{まもる}} — \emph{mamoru} — \textquotedblleft{}to protect; to guard\textquotedblright{}.

\begin{quote}
\textbf{\ja{さむらいは} \ja{おうさまを} \ja{まもっていました。}} — \emph{samurai wa ōsama o mamotte imashita} — \textquotedblleft{}The samurai was guarding the king.\textquotedblright{}
\end{quote}

Polite \textbf{\ja{まもります}}, \emph{mamorimasu}; a godan verb. Past (polite): \textbf{\ja{さむらいが} \ja{むらを} \ja{まもりました。}}, \emph{samurai ga mura o mamorimashita}.

\begin{cousinweb}
\textbf{\ja{まもる}} (\emph{mamoru}) is a native Japanese word (wago), inherited rather than borrowed from Chinese; its kanji was attached later as a kun'yomi (JMdict).
\end{cousinweb}

\begin{grammarlens}[title={What you've built}]
One more word for telling what happened.
\end{grammarlens}

\subsection*{Your turn}
\begin{itemize}
\raggedright
  \item \emph{Say it:} \emph{mamoru}
  \item \emph{Say it:} \emph{mamoru}, once more
  \item \emph{Say it:} \emph{samurai wa ōsama o mamotte imashita} — The samurai was guarding the king.
  \item \emph{Recall:} say the Japanese for a princess, then the Japanese for a samurai, then say \emph{mamoru} again
\end{itemize}

\begin{tcolorbox}[breakable,colback=teal!4,colframe=teal!35!black,title={Before you move on}]
What does \ja{まもる} mean? (\textquotedblleft{}To protect; to guard\textquotedblright{}.) Say it once more.
\end{tcolorbox}
\section[\texorpdfstring{\ja{とのさま} (tonosama)}{とのさま (tonosama)}]{\ja{とのさま} (tonosama) — a feudal lord}

\label{lesson:JA-C299-tonosama}



\begin{quote}
Before the new one: say the Japanese for a samurai, then the Japanese for to protect.
\end{quote}

\subsection*{You'll want to know: \ja{とのさま}}
\textbf{\ja{とのさま}} — \emph{tonosama} — \textquotedblleft{}a feudal lord\textquotedblright{}.

\begin{quote}
\textbf{\ja{とのさまは} \ja{さむらいを} \ja{よびました。}} — \emph{tonosama wa samurai o yobimashita} — \textquotedblleft{}The lord sent for the samurai.\textquotedblright{}
\end{quote}

\begin{cousinweb}
\textbf{\ja{とのさま}} (\emph{tonosama}) is a native Japanese word (wago), inherited rather than borrowed from Chinese, built from \emph{tono}, \textquotedblleft{}lord\textquotedblright{}, and \emph{sama}, \textquotedblleft{}the honorific\textquotedblright{} (JMdict, KANJIDIC2).
\end{cousinweb}

\begin{grammarlens}[title={What you've built}]
One more word for telling what happened.
\end{grammarlens}

\subsection*{Your turn}
\begin{itemize}
\raggedright
  \item \emph{Say it:} \emph{tonosama}
  \item \emph{Say it:} \emph{tonosama}, once more
  \item \emph{Say it:} \emph{tonosama wa samurai o yobimashita} — The lord sent for the samurai.
  \item \emph{Recall:} say the Japanese for a samurai, then the Japanese for to protect, then say \emph{tonosama} again
\end{itemize}

\begin{tcolorbox}[breakable,colback=teal!4,colframe=teal!35!black,title={Before you move on}]
What does \ja{とのさま} mean? (\textquotedblleft{}A feudal lord\textquotedblright{}.) Say it once more.
\end{tcolorbox}
